% Autogenerated translation of index.md by Texpad
% To stop this file being overwritten during the typeset process, please move or remove this header

\documentclass[12pt]{book}
\usepackage{graphicx}
\usepackage{fontspec}
\usepackage[utf8]{inputenc}
\usepackage[a4paper,left=.5in,right=.5in,top=.3in,bottom=0.3in]{geometry}
\setlength\parindent{0pt}
\setlength{\parskip}{\baselineskip}
\setmainfont{Helvetica Neue}
\usepackage{hyperref}
\pagestyle{plain}
\begin{document}

\hrule
layout: default

\section*{title: Homepage}

\section*{About}

\begin{verbatim}
<div class="about-grid" markdown="1">
<div class="about-copy" markdown="1">

I am a fourth-year PhD student in the Department of Mathematics at Brown University. My advisor is [Benoît Pausader](https://www.math.brown.edu/bpausade/).

I completed my undergraduate studies at the University of Science and Technology of China (USTC). You can find my CV [here](CV_hwr.pdf).

I am interested in analysis and partial differential equations. My current research focuses on nonlinear dispersive PDEs and kinetic equations, especially long-time behavior and stability problems.

**Contact:** wenrui_huangAtbrownDotedu

</div>
<figure class="about-portrait">
  <img src="wenrui2.png" alt="Portrait of Wenrui Huang">
</figure>
</div>

\end{verbatim}

\section*{Publications \& Preprints}

\begin{enumerate}
\item \href{https://arxiv.org/abs/2608.04402}{Long-Time Dynamics and Modified Scattering for the Coulomb Vlasov–Hartree System}, (with M. Xie), arXiv:2608.04402
\texttt{$<$details class="paper-abstract"$>$}
 \texttt{$<$summary$>$}\texttt{$<$span class="visually-hidden"$>$}Toggle abstract\texttt{$<$/span$>$}\texttt{$<$/summary$>$}
 \texttt{$<$p$>$}We study the three-dimensional Coulomb Vlasov-Hartree system, a classical-quantum model for Bose-Fermi mixtures. For small, regular, and localized data, the paper proves global well-posedness and identifies the coupled long-time dynamics: logarithmic corrections appear in both the Vlasov characteristics and the Hartree phase.\texttt{$<$/p$>$}
\texttt{$<$/details$>$}
\item \href{https://arxiv.org/abs/2602.21344}{Scattering map for the Vlasov-Poisson system with a harmonic repulsive potential}, (with H. Kwon), arXiv:2602.21344
\texttt{$<$details class="paper-abstract"$>$}
 \texttt{$<$summary$>$}\texttt{$<$span class="visually-hidden"$>$}Toggle abstract\texttt{$<$/span$>$}\texttt{$<$/summary$>$}
 \texttt{$<$p$>$}We study scattering for the Vlasov-Poisson system with a repulsive harmonic potential in dimensions at least two. The work proves modified scattering and constructs wave operators through a lens-transform viewpoint, with weaker assumptions on the initial data than in some earlier treatments.\texttt{$<$/p$>$}
\texttt{$<$/details$>$}
\item \href{https://arxiv.org/abs/2601.10030}{Stability and instability of small BGK waves}, (with D. Bian, E. Grenier and B. Pausader), arXiv:2601.10030
\texttt{$<$details class="paper-abstract"$>$}
 \texttt{$<$summary$>$}\texttt{$<$span class="visually-hidden"$>$}Toggle abstract\texttt{$<$/span$>$}\texttt{$<$/summary$>$}
 \texttt{$<$p$>$}This paper analyzes the linear behavior of small Bernstein-Green-Kruskal waves. The stability or instability mechanism is tied to the sign of the derivative of the energy distribution at zero energy.\texttt{$<$/p$>$}
\texttt{$<$/details$>$}
\item \href{https://arxiv.org/abs/2412.13434}{The Vlasov–Poisson system with a perfectly conducting wall: Convex domains}, (with B. Pausader and M. Suzuki), to appear in \textbf{Comm. Math. Phys.}
\texttt{$<$details class="paper-abstract"$>$}
 \texttt{$<$summary$>$}\texttt{$<$span class="visually-hidden"$>$}Toggle abstract\texttt{$<$/span$>$}\texttt{$<$/summary$>$}
 \texttt{$<$p$>$}We investigate the Vlasov-Poisson system in convex domains with a perfectly conducting boundary. Under localization and geometric assumptions, the asymptotic particle velocities are constrained by the asymptotic domain, and the dynamics exhibit a modified scattering behavior.\texttt{$<$/p$>$}
\texttt{$<$/details$>$}
\item \href{https://arxiv.org/abs/2411.02737}{Modified wave operators for the Hartree equation with repulsive Coulomb potential}, arXiv:2411.02737
\texttt{$<$details class="paper-abstract"$>$}
 \texttt{$<$summary$>$}\texttt{$<$span class="visually-hidden"$>$}Toggle abstract\texttt{$<$/span$>$}\texttt{$<$/summary$>$}
 \texttt{$<$p$>$}This work treats a final-state problem for the Hartree equation with repulsive Coulomb interaction. For small and sufficiently localized asymptotic data, it constructs a unique global solution with the prescribed modified scattering profile.\texttt{$<$/p$>$}
\texttt{$<$/details$>$}
\item \href{https://link.springer.com/article/10.1007/s00023-026-01735-7}{Scattering of the Vlasov–Riesz system in three dimensions}, (with H. Kwon), \textbf{Annales Henri Poincaré} (2026), \href{https://arxiv.org/abs/2407.16919}{arXiv:2407.16919}.
\texttt{$<$details class="paper-abstract"$>$}
 \texttt{$<$summary$>$}\texttt{$<$span class="visually-hidden"$>$}Toggle abstract\texttt{$<$/span$>$}\texttt{$<$/summary$>$}
 \texttt{$<$p$>$}We study long-time behavior for the three-dimensional Vlasov-Riesz system. The paper establishes scattering for small data in sub-Coulomb regimes, modified scattering near the Coulomb case, and corresponding wave operators, including polynomial corrections in the asymptotic description.\texttt{$<$/p$>$}
\texttt{$<$/details$>$}
\end{enumerate}

\section*{Teaching}

\subsection*{Brown University}

\begin{itemize}
\item Fall 2025 — TA — MATH 0190 Calculus II (Physics/Engineering)
\item Fall 2023, Spring 2024, Fall 2026 — TA — MATH 0180 Calculus III
\item Spring 2023 — TA — MATH 0100 Calculus II
\end{itemize}

\subsection*{University of Science and Technology of China}

\begin{itemize}
\item Spring 2022 — TA — Riemannian Geometry
\item Spring 2021, Spring 2022 — TA — Topology
\item Fall 2021 — TA — Differential Geometry
\end{itemize}

\section*{Talks}

\begin{itemize}
\item GLESPA seminar, Brown University, April 2026.
\item Poster session in Workshop on Kinetic Theory and Fluids, University of Wisconsin-Madison, March 2025.
\item GLESPA seminar, Brown University, April 2024.
\item Brown-Yale PDE seminar,  April 2024.
\item Brown-Yale PDE seminar, November 2023.
\end{itemize}

\end{document}
